%!TEX program = uplatex
\RequirePackage{plautopatch}
\documentclass[uplatex,dvipdfmx]{jlreq}
\usepackage{amsmath,amssymb,amscd,amsthm}

\pagestyle{empty}

\begin{document}

%======================================================================
% 表紙・概要：栗原将人「岩澤数学への招待」
%======================================================================
\begin{center}
  {\Large\bfseries 岩澤数学への招待}

  \vspace{1pc}
  {\large 概要}
\end{center}
\begin{flushright}
  栗原　将人（都立大・理）
\end{flushright}

今回の企画のテーマは岩澤数学です。
一昨年お亡くなりになられた岩澤健吉先生は、
さまざまな独創的理論をお創りになられましたが、
その中でも有名な岩澤分解と岩澤理論について、
その入門の入門をしたいと考えております。
今回の企画の趣旨は、他分野の方々や多くの学生さん達を対象に考えた
本当の入門の入門でして、
知られているさまざまな結果をいろいろと述べるとか、
最先端の結果を述べるとかいうことを目標とはしていません。
聴衆として実際に想定しておりますのは、
整数論を専門に研究されている方々ではなく、
だいたい学部の３年生くらいの知識を持った方々ですので、
どうぞ興味のある方は御参集下さい。

岩澤分解は、Lie 群の基本的な定理です。
Hilbert の第５問題や、($L$)-群の理論などともからめて、
佐武先生に講義していただく予定です。

整数論における、いわゆる岩澤理論は、
円分体のイデアル類群に関して、美しい関係が成立すること
（岩澤類数公式）の発見から始まりました（[1]）。
尾崎さんには、
整数論の重要な研究対象の一つであるイデアル類群を考えることの意味、
おもしろさから話し始めてもらい、
岩澤類数公式、及びそこに現れる岩澤不変量についての解説をしていただきます。
また、岩澤不変量を考えるおもしろさについても語っていただく予定です。

岩澤理論の主要な核となる Iwasawa Main Conjecture は
$p$ 進解析的なものと $p$ 進代数的なものとの結びつき、という形で表現されます。
土曜日の前半は、私が $p$ 進ゼータを中心にして、$p$ 進解析的な側面から話します。
Bernoulli 数の $p$ 進的性質等、実例を中心として話し、予備知識は特に何も仮定しません。
講演全般、特に土曜の講演に対しての Introduction も行います。

土曜の後半はまず、市村先生に Iwasawa Main Conjecture へ至る道について、
すなわち $p$ 進解析的なものと $p$ 進代数的なものがいかに結びつくのか、
ということを、わかりやすく説明していただきます。

岩澤理論はたぶん最初に考えられていたより、ずっと適用範囲が広い理論であり、
そのアイディアがたくさんのことに応用されています。
このことは、たとえばワイルスによるフェルマー予想の証明で有名になりました。
現在、岩澤理論は整数論の一支脈をなし、さまざまな研究がなされています。
最後の時間で、加藤先生に岩澤理論に関する夢について語っていただこうと思っております。

なお、岩澤数学についての一般向けの解説書としては、[2] があります。

\vspace{1pc}

\begin{center}
  \textsc{参考文献}
\end{center}
\begin{enumerate}
  \item K.~Iwasawa,
        On $\Gamma$-extensions of algebraic number fields,
        \textit{Bull.~Amer.~Math.~Soc.}~\textbf{65} (1959), 183--226.
  \item 数学のたのしみ 15 号、日本評論社（1999年10月）
\end{enumerate}

\newpage

%======================================================================
% 講演１：佐武一郎「岩澤分解について」
%======================================================================
\begin{center}
  {\Large\bfseries 岩澤分解について}
\end{center}

\bigskip
\hspace*{\fill}{\large\bfseries 佐武　一郎（東北大 / UC Berkeley）}

\bigskip

岩澤分解はリー群論で基本的な、非常に有用な定理の一つです。
線形代数でよく知られたシュミットの直交化法を行列を使って表すと、
丁度 $GL(n,\mathbf{R})$ の場合の岩澤分解が得られます。
こういう具体例から始めて、最初、半単純群、代数群の場合の一般定理を説明しようと思います。

この定理は岩澤先生の有名な論文 [I49] の中で、$L$-群の構造定理のための
補助定理として証明されたものです。
講演の後半で、この構造定理、さらに $L$-群に関する岩澤予想、
それとヒルベルトの第５問題との関係等についても話したいと思います。
--- 後年の岩澤理論の考え方の萌芽が、すでにこの論文の中にあることを
感知して頂ければ幸いです。

\vspace{6pt}

\begin{center}
  \textsc{参考文献}
\end{center}
\begin{enumerate}
  \item[{[I49]}] K.~Iwasawa,
        On some types of topological groups,
        \textit{Ann.~of Math.}~\textbf{50} (1949), 507--557.
  \item[{[S]}] 佐武一郎,
        岩澤先生と位相群論,
        数学のたのしみ, no.~15 (1999年10月), 9--21.
\end{enumerate}

\newpage

%======================================================================
% 講演２：尾崎学「岩澤類数公式と岩澤不変量のおもしろさ」
%======================================================================
\begin{center}
  {\Large\bfseries 岩澤類数公式と岩澤不変量のおもしろさ}
\end{center}

\bigskip
\hspace*{\fill}{\large\bfseries 尾崎　学（島根大・総合理工）}

\bigskip

有理数体 $\mathbb{Q}$ が有理整数環 $\mathbb{Z}$ を部分環として持つように、
代数体 $k$（$\mathbb{Q}$ の有限次拡大体）も普通の整数（有理整数）の一般化である
代数的整数の環 $\mathcal{O}_k$ という部分環を持ちます。
古代ギリシア以来知られているように $\mathbb{Z}$ に於いては素因数分解定理が成り立っています。
つまり、どんな $0$ でない有理整数もいくつかの素数の積として本質的に一通りの方法で表されます。
しかし、幸か不幸か $\mathcal{O}_k$ では一般にこの素因数分解定理が成立しないのです。
歴史的にこの困難との遭遇とその克服は、19世紀に Fermat の最終定理を巡って起きました。
講演ではその辺から始めて、代数的整数論の主要な研究主題であって、
$\mathcal{O}_k$ での素因数分解の一意性の破綻の様子を表す「イデアル類群」というものについて、
例も交えてまず説明します。

次に固定された素数 $p$ に対して円の $p$ 冪分体達
$\mathbb{Q}(\zeta_{p^n})$（$n \geq 1$）を考えます。
すると、これらの円分体達の類数（＝イデアル類群の位数）の間には
岩澤類数公式と呼ばれる美しい関係式が成り立ちます。
代数体の類数というものは計算することすら困難で、類数の「分かりやすい」公式は
特別な場合を除いて知られていません。
それにもかかわらず、このような関係式が存在するのは非常に驚くべきことです。
岩澤類数公式に現れる岩澤不変量には円の $p$ 冪分体達のイデアル類群の情報が込められています。
これら岩澤類数公式と岩澤不変量について、実例や予想・問題も交えて解説します。

\vspace{1pc}

\begin{center}
  \textsc{参考文献}
\end{center}
\begin{enumerate}
  \item K.~Iwasawa,
        On $\Gamma$-extensions of algebraic number fields,
        \textit{Bull.~Amer.~Math.~Soc.}~\textbf{65} (1959), 183--226.
  \item L.~C.~Washington,
        \textit{Introduction to Cyclotomic Fields} (2nd edition),
        Graduate Texts in Mathematics \textbf{83},
        Springer-Verlag, 1996.
  \item 黒川信重・栗原将人・斎藤毅:
        岩波講座 現代数学の基礎 数論３
\end{enumerate}

\newpage

%======================================================================
% 講演３：栗原将人「$p$ 進ゼータのたのしみ」
%======================================================================
\begin{center}
  {\Large\bfseries $p$ 進ゼータのたのしみ}
\end{center}

\bigskip
\hspace*{\fill}{\large\bfseries 栗原　将人（都立大・理）}

\bigskip

\noindent\textbf{１．}
Riemann zeta
\[
  \zeta(s) = \sum_{n=1}^{\infty} n^{-s}
\]
の値に関しては、Euler 以来
$\zeta(2) = \pi^2/6$、$\zeta(4) = \pi^4/90$、$\zeta(6) = \pi^6/945$ \ldots
などが知られているが、この有理数部分を取り出してみると、
そこには $p$ 進的なつながりが存在する。
このことをさまざまな実例を通じて見ていこうと思う。
たとえば、Kummer の合同式として知られている有名な合同式が
このような $p$ 進的性質の例であるが、
ここでは Kummer の合同式以外の例もいろいろと見ていきたい。
そして、それらが $p$ 進 zeta の存在から説明できること、
特に $p$ 進 zeta が「岩澤関数」であることを述べたいと思う。
さらに未解決の問題もいくつか述べる
（このような単純な部分にも未解決問題が存在する）。
普段、$\mathbf{R}$ や $\mathbf{C}$ の世界に住んでいる方々にも
$p$ 進世界の一端を紹介したい。

\medskip

\noindent\textbf{２．}
こうして $p$ 進ゼータの背後には、ある種の冪級数もしくは多項式があることがわかるのであるが、
この冪級数、もしくは多項式の性質をいろいろと調べていく。
この多項式は純粋に $p$ 進解析的に作られたにもかかわらず、
整数論的な情報を驚くほど多く含んでいるのである。

\vspace{1pc}

\begin{center}
  \textsc{参考文献}
\end{center}
\begin{enumerate}
  \item K.~Iwasawa,
        \textit{Lectures on $p$-adic $L$-functions},
        Annals of Mathematics Studies \textbf{74},
        Princeton University Press, 1972.
\end{enumerate}

\newpage

%======================================================================
% 講演４：市村文男「岩澤主予想への道」
%======================================================================
\begin{center}
  {\Large\bfseries 岩澤主予想への道\\
  \normalsize --- Stickelberger の定理とその周辺}
\end{center}

\bigskip
\hspace*{\fill}{\large\bfseries 市村　文男（横浜市大・理）}

\bigskip

円分体 $K_n = \mathbb{Q}(e^{2\pi i/p^{n+1}})$ の類数の
plus part $h_n^+$、odd part $h_n^-$ についての「解析的類数公式」というものがあります。
19世紀に Kummer によって得られたこの不思議な公式を代数的に理解したいというのが
岩澤主予想の出発点です。
$h_n^-$ についての公式は、いくつかの Bernoulli 数を用いて表わされますが、
この Bernoulli 数の「親玉」が Stickelberger 元といわれる、群環
$\mathbb{Z}[\mathrm{Gal}(K_n/\mathbb{Q})]$ の元です --- $\theta_n$ とかく ---。
栗原さんの講演で述べられるように、Bernoulli 数を $p$ 進的につないだものが
$p$ 進 $L$ 関数です。
これらをもとにして、岩澤先生は
「$p$ 進 $L$ 関数 $=$ Stickelberger 元（達）」という等式を得ました。
Stickelberger 元は、代数的にも $K_n$ の類群 $\mathrm{Cl}_n$ と関係しています。
すなわち、$\theta_n$ の自然な作用で $\mathrm{Cl}_n$ は消えてしまうのです
（Stickelberger の定理）。
以上の事から、$\theta_n$ は、従って、解析的につくられた $p$ 進 $L$ 関数が
$\mathrm{Cl}_n$ の odd part $\mathrm{Cl}_n^-$ の代数構造を
（ほぼ）完全に知っているのではないかと期待したくなります。
これが岩澤主予想です。以上のような事をじっくり話したいと思っています。

\newpage

%======================================================================
% 講演５：加藤和也「岩澤理論とその夢」
%======================================================================
\begin{center}
  {\Large\bfseries 岩澤理論とその夢}
\end{center}

\bigskip
\hspace*{\fill}{\large\bfseries 加藤　和也（東大・数理）}

\bigskip

岩澤理論は、
\begin{center}
  ゼータ関数
  \raisebox{-5pt}{$\overleftrightarrow{\text{\scriptsize{~~ふしぎな関係~~}}}$}
  イデアル類群
\end{center}
というふうに、解析的な存在であるゼータ関数と代数的な存在であるイデアル類群という、
異類のものの間に存在するふしぎな関係を探る理論である。

この岩澤理論は、
\begin{center}
  ゼータ関数
  \raisebox{-5pt}{$\overleftrightarrow{\text{\scriptsize{~~ふしぎな関係~~}}}$}
  楕円曲線の有理点の群
\end{center}
（Coates と Wiles がはじめた）のように、
\begin{center}
  ゼータ関数
  \raisebox{-5pt}{$\overleftrightarrow{\text{\scriptsize{~~ふしぎな関係~~}}}$}
  重要な代数的な群
\end{center}
の形に一般化されつつある。
Wiles によるフェルマーの最終定理の証明、谷山--志村予想の大きな部分の証明も、
この一般化の流れの中にある。

こうしたことについて解説したい。

\end{document}